\section{Materials and methods}

\subsection{Implementation, validation, and definition of
success}\label{sec:implementation-en}

For a base \(a\) coprime to the number \(N\), the sequence \(a^{x} \bmod N\) is
periodic; its order \(r\) is the smallest positive integer satisfying
\(a^{r} \equiv 1 \pmod{N}\). The quantum part estimates it by phase estimation and
inverse QFT; the classical part derives and verifies the factors through rational
approximation and greatest common divisors.

The main instance is \(N = 15\), \(a = 7\), with eight control and four work qubits. Its
order is \(r = 4\), and the ideal control-register outcomes are 0, 64, 128, and 192. Each
measured execution, or \emph{shot}, returns one outcome; multiple shots form the
post-processing histogram. Success means recovering non-trivial factors, not reconstructing
the exact order. Three ideal peaks yield the factors, for 75\% overall probability.
Classical verification is permissive, however: it accepts 63 of 256 values, so uniform
outcomes still succeed with probability \(63/256 \simeq 24.61\%\). This instance- and
post-processing-specific reference is not quantum-state fidelity; it shows that correct
factorizations can survive after the distribution loses much of its structure
\cite{en-smolin2013}.

Circuit correctness is validated independently of factorization success by checking
modular operators, restored ancillas, and theoretical distributions. The ideally
validated instances \(N = 21\), \(a = 2\) and \(N = 35\), \(a = 6\) use Beauregard's
reversible arithmetic \cite{en-beauregard2002}; only \(N = 21\) enters the exploratory
noisy AQFT tests. Qiskit/Aer, scikit-learn, Stim, and PyMatching are used, recording versions,
parameters, random seeds, and circuit identifiers.

\subsection{Post-processing, classification, and mitigation}

Three strategies share the same factor extractor, tie-breaking rule, and 1,024-shot
histograms, one per quantum iteration. TOP-1 (Method 1) checks only the most frequent
outcome; TOP-4 checks up to four frequency-ranked candidates; Method 2 applies TOP-4 only
when a classifier predicts success for the dominant candidate. Comparing Method 2 with
unfiltered TOP-4 ablates the classifier and separates its contribution from that of
testing more candidates.

Two illustrative scenarios, reference UC1 and stress case UC2, include depolarization,
thermal relaxation, and readout errors. For each, 2,000 histograms with parameters varied
within \(\pm50\%\) of nominal are split 60\%/20\%/20\% into training, selection, and test
sets. Random forests, support vector machines (SVMs), and multilayer perceptrons are
compared on F1; the selected model is tested on 400 independent histograms per scenario.
The primary metric is iterations to obtain the factors, not classifier accuracy or
runtime. Thirty paired replicates per scenario run for at most 50 iterations; failures
receive the conventional value 51. The analysis uses Wilcoxon--Pratt tests with Holm
correction. Two- and three-level zero-noise extrapolation (ZNE) scales only depolarizing
parameters in simulation \cite{en-temme2017,en-takagi2023}.

\subsection{Quantum error correction and syndrome decoding}

Quantum error correction distributes one logical qubit across several physical qubits.
Parity checks, called syndromes, reveal errors without reading the logical information;
a decoder selects the correction. Repetition, Steane \([[7,1,3]]\), and surface codes
\cite{en-steane1996,en-fowler2012,en-googleQECBelowThreshold2025} are studied in memory circuits through the logical
failure probability \(p_{L}\), the probability that information remains wrong after
decoding.

The rotated surface code uses distances \(d = 3,5,7,9\) (the number of errors the code can
tolerate), \(X\)- and \(Z\)-basis memories, and \(d\) syndrome rounds, with noise on
gates, preparation, and measurement. On the same detection events, reference
minimum-weight perfect matching (MWPM) is compared with a stand-alone neural network and
a hybrid decoder. The hybrid learns when to flip a likely-wrong matching decision, with
an intervention threshold chosen in validation and fixed for testing. Crosstalk between
adjacent qubits is added to nominal noise; model calibration, network input size, and the
analytical BP+OSD decoder (belief propagation and ordered statistics decoding) are also
examined \cite{en-alphaqubit2024,en-roffe2020}.

\subsection{Sensitivity of Shor's algorithm to gate errors}

A separate experiment applies independent Pauli errors after compiled gates, excluding
virtual RZ rotations, using 20 replicates of 4,096 shots per value. The total non-identity
error probability \(p_{g}\) is a per-gate proxy; without complete fault-tolerant
compilation, it cannot be converted into the memory circuits' logical failure
probability \(p_{L}\).

\subsection{Optimization through the approximate QFT}

The AQFT removes small-angle controlled rotations, trading precision for fewer gates,
lower depth, and fewer opportunities for noise \cite{en-barencoAQFT1996}.
Truncation of Shor's final inverse QFT
for \(N = 15\) is distinguished from truncation of transforms inside the modular
arithmetic for \(N = 21\). Separate phase-estimation benchmarks use six-, eight-, and
ten-qubit registers with three phases per size. The noise model uses gate infidelities
and readout errors from an offline FakeSherbrooke calibration, without additional
\(T_{1}/T_{2}\) contributions and with common layouts and random seeds. The truncation
degree is chosen on selection samples and assessed on disjoint test samples (in the
\(N = 15\) pilot, 8,192 shots per configuration, split equally).
